% critique.tex — f172v, vv. 835–876 : commentaire critique.
% Parties textuelles du commentaire de l'auteur (travaux antérieurs, Harmattan 2023), texte et notes
% conservés ; répartition, harmonisation et passages signalés : repartition.md.
\cc{839}{ὦ πόλεις.}{Le texte du manuscrit de \italique{Ravenne} comporte au vers 839 une adresse à la Cité :\citComm{τῷ γὰρ εἰκὸς ὦ πόλεις (sic.)}La vraisemblable quasi-coïncidence des dates entre la mort d'Eupolis et la représentation des \italique{Thesmophories}\nc{Voir l'introduction et la question de la date. Avec Rogers, je penche pour l'hypothèse que les \italique{Thesmophories} furent représentées aux Dionysies de 410. Si j'ai vu juste avec Natalia, Eupolis serait donc mort l'été précédent la première représentation de la pièce.} m'incite à penser qu'il y a là plutôt une évocation d'Eupolis, auteur intarissable sur la mère d'Hyperbolos, et à laquelle il a même consacré une pièce entière : son \italique{Maricas} (voir le commentaire du vers 840).
Sur la mort d'Eupolis, voir le commentaire de la pièce.}
\cc{853}{πικρὰν ἑλένη.}{Quand Kritylla prononce `πικρὰν ἑλένην ὄψει τάχ᾽', elle ne songe pas un seul instant à la mythique Hélène, mise nouvellement en scène par Euripide avec réécriture du scénario. C'est Hésychius encore une fois – décidément le plus utile des lexicographes pour les \italique{Thesmophories}~! – qui nous le confirme (ε 1995)~:\citComm{ἑλένη· λαμπάς, δετή\\
\francais{hélène~: flambeau, torche.}}}
